% page 19 du fac-simile (image p-018 du scan)
% bloc 160.0 x 224.7 mm ; marge gauche 42.0 mm
\pgc{20.320mm}{0.000mm}{
\l{\cel{55}11}
\l{}
\l{}
\l{\sou{agnoskar}. (\sou{trans}.) Lasar enirar sua spirito kom vera, valida,}
\l{\cel{3}yusta, legale konstitucita. - L.}
\l{}
\l{\sou{agoniar}. (\sou{netrans.}) Esar la objekto di la lukto lasta da naturo}
\l{\cel{3}kontre morto, en la korpo di vivanto. - DEFIRS.}
\l{}
\l{\sou{agorafobio}. Pavoro morba, da ula personi, kande li devas trans-}
\l{\cel{3}irar placo, strado, ponto. - EFIRS.}
\l{}
\l{\sou{agosto}. La monato okesma di la yaro kalendarala.}
\l{}
\l{\sou{agro}.\cel{2}I.Tereno sen-tekta e plata, limitizita por uzeso parti-}
\l{\cel{3}kulara. - II. Tereno por kultivado, ne cirkondata da muri.}
\l{\cel{2}- DEFIRS.}
\l{}
\l{\sou{agrafo}. I. (\sou{arkitekt.}) Krampono fera qua inter-adherigas la}
\l{\cel{3}quadri di muro. - II. (\sou{Vesto}). Roketo metala qua uzesas}
\l{\cel{3}por interjuntar la du bordi inter-opozanta di vesto, fixig-}
\l{\cel{3}ita sur ica, ed eniranta ringo, fixigita sur ita. - DEFIS.}
\l{}
\l{\sou{agreabla}.\cel{2}(ad).\cel{2}Qua donas plezuro (ad ulu). - EFIS.}
\l{}
\l{\sou{agregar}.\cel{2}(\sou{trans.}) I. Interunionar - por formacar ek li en-}
\l{\cel{3}semblo - molekuli, partikuli materio. - II. Unionar ad}
\l{\cel{3}ensemblo, a korpo. - III. (\sou{en la universitato di Francia})}
\l{\cel{3}\sou{agregito}\cel{1}: la persono qua, pro sucesir lor la konkurso pri}
\l{\cel{3}Io, deklaresas kapabla pri docar en liceo, en fakultato.}
\l{\cel{3}(Ta grado esas min alta kam doktoreso). - DEFIRS.}
\l{}
\l{\sou{agrimonio}.\cel{2}(\sou{bot}.) Planto rozacea, kun mikra flori qui esas}
\l{\cel{3}dispozita spike, olim uzita kom astriktivo. - L.\cel{1}\sou{agrimonia}}
\l{\cel{3}\sou{eupatoria}.- DEFIS.}
\l{}
\l{\sou{agronomo}.\cel{2}Persono experta pri la cienco \textquotedbl{}agrokultivo\textquotedbl{}.-DEFIRS.}
\l{}
\l{\sou{agronomio}. La cienco qua donas la reguli di agrokultivo.}
\l{\cel{3}- DEFIRS .}
\l{}
\l{\sou{agulo}. I. Stangeto ek stalo polisita, diqua un extremajo esas}
\l{\cel{3}pinta, e la altra, obtuzigita, ed un-trua, por recevar filo.}
\l{\cel{3}- II. Mikra stango di qua un extremajo esas pinta. (\sou{Agulo}}
\l{\cel{3}\sou{di relvoyo}\cel{1}: parti di reli movebla, uzata por igar la treni}
\l{\cel{3}pasar bifurke de ca a ta relvoyo). - FIS.}
\l{}
\l{\sou{ah !}\cel{2}Interjeciono qua indikas ke onu quik expresos emoco}
\l{\cel{3}forta qua kaptis la psiko (joyo, doloro, iraco, e c.)}
\l{}
\l{\sou{ailanto}. (\sou{bot.})\cel{2}Planto ek la familio \textquotedbl{}nutacei\textquotedbl{}. -}
\l{\cel{3}L.\cel{1}\sou{Ailanthus}.}
}
